\section{Rotten Oranges}
\label{sec:graph-rotten-oranges}

\problemstatement{Given a grid where each cell is $0$ (empty), $1$
(fresh orange), or $2$ (rotten orange), every minute any fresh orange
\textbf{4-directionally adjacent} to a rotten orange becomes rotten.
Find the minimum number of minutes until no fresh orange remains, or
$-1$ if some fresh orange can never be reached.}

\begin{patternbox}
\emph{``Multiple sources spread outward simultaneously, find the time
for the spread to finish''} is solved by \textbf{multi-source BFS}: push
\textbf{every} initially rotten orange into the queue at the start
(not just one), then run ordinary BFS -- the key realization is that
BFS already processes nodes in order of distance from \emph{whichever}
source is nearest, automatically, when all sources start in the queue
together at time $0$.
\end{patternbox}

\subsection*{Why starting from all sources at once is correct}

If BFS were run separately from each rotten orange one at a time, the
correctness would still hold but the bookkeeping (combining several
separate distance computations, taking a minimum over all of them)
would be needlessly complex. Pushing every source into the \emph{same}
queue at the very start works because BFS's queue already guarantees
nodes are processed in non-decreasing distance order \emph{from
whichever source reaches them first} -- exactly the quantity this
problem wants (a fresh orange rots at the minute the \emph{nearest}
rotten orange's spread reaches it).

\subsection*{Algorithm}

\begin{enumerate}
  \item Scan the grid: push every rotten orange's coordinates into a
        queue, and count the total number of fresh oranges.
  \item Run BFS level by level (processing the entire current queue
        before moving to the next level represents one minute
        passing): for each rotten orange popped, check its $4$
        neighbors; any fresh neighbor becomes rotten, decrement the
        fresh count, and push it for the next level.
  \item Track the number of levels (minutes) processed. After BFS
        completes, if any fresh oranges remain uncounted, return $-1$;
        otherwise return the number of minutes elapsed.
\end{enumerate}

\subsection*{C++ implementation}

\begin{lstlisting}
#include <bits/stdc++.h>
using namespace std;

int orangesRotting(vector<vector<int>>& grid) {
    int rows = grid.size(), cols = grid[0].size();
    queue<pair<int,int>> q;
    int freshCount = 0;

    for (int r = 0; r < rows; r++) {
        for (int c = 0; c < cols; c++) {
            if (grid[r][c] == 2) q.push({r, c});
            else if (grid[r][c] == 1) freshCount++;
        }
    }

    int minutes = 0;
    int dr[] = {-1, 1, 0, 0};
    int dc[] = {0, 0, -1, 1};

    while (!q.empty() && freshCount > 0) {
        int levelSize = q.size();
        bool rottedThisLevel = false;

        for (int i = 0; i < levelSize; i++) {
            auto [r, c] = q.front();
            q.pop();

            for (int dir = 0; dir < 4; dir++) {
                int nr = r + dr[dir], nc = c + dc[dir];
                if (nr >= 0 && nr < rows && nc >= 0 && nc < cols &&
                    grid[nr][nc] == 1) {
                    grid[nr][nc] = 2;
                    freshCount--;
                    q.push({nr, nc});
                    rottedThisLevel = true;
                }
            }
        }
        if (rottedThisLevel) minutes++;
    }
    return freshCount == 0 ? minutes : -1;
}

int main() {
    vector<vector<int>> grid = {
        {2, 1, 1},
        {1, 1, 0},
        {0, 1, 1}
    };
    cout << orangesRotting(grid) << endl; // 4
    return 0;
}
\end{lstlisting}

\subsection*{Dry run}

\dryrun{Take the $3\times3$ grid above (one rotten orange at
$(0,0)$).}

Initial queue: $[(0,0)]$, $\textit{freshCount}=6$. Minute $1$: pop
$(0,0)$, rot neighbors $(1,0)$ and $(0,1)$ (both fresh); push both;
$\textit{freshCount}=4$. Minute $2$: pop $(1,0)$, rot $(1,1)$; pop
$(0,1)$, rot $(0,2)$ ($(1,1)$ is already rotten by this point, so
$(0,1)$'s other neighbor is skipped); $\textit{freshCount}=2$. Minute
$3$: pop $(1,1)$, rot $(2,1)$; pop $(0,2)$, no fresh neighbors;
$\textit{freshCount}=1$. Minute $4$: pop $(2,1)$, rot $(2,2)$;
$\textit{freshCount}=0$. The queue still holds $(2,2)$, but the loop
exits immediately since $\textit{freshCount}=0$ -- final answer: $4$
minutes.

\begin{complexitybox}
\textbf{Time Complexity.} $O(\textit{rows} \times \textit{cols})$ --
every cell is pushed onto the queue and processed at most once.
\textbf{Space Complexity.} $O(\textit{rows} \times \textit{cols})$ for
the queue in the worst case (a grid that is entirely rotten oranges
initially).
\end{complexitybox}

\begin{takeawaybox}
Multi-source BFS -- seeding the queue with every source simultaneously
rather than running BFS once per source -- is the standard technique
whenever a problem describes several things spreading or growing
outward at the same rate and asks for the time until the spread
finishes or the distance to the nearest source. The very next section,
01 Matrix, applies this exact same technique to a differently-phrased
problem.
\end{takeawaybox}
